%% ethics.tex

\section{Ethics Statement and Design Recommendations}
\label{sec:ethics}

\subsection{System Identity and Authorship}

The system studied in this paper, \system{}, is a multi-agent LLM operating system developed and maintained by the authors of this paper. The attack vector studied (V2: \code{TopologyMutationStrategy} privilege escalation via missing authorization on the role-replacement primitive) was identified during a security review of the authors' own codebase, not through adversarial probing of a third-party system. We are the maintainers of \system{}; the work is therefore a \emph{self-study of an authorization flaw in our own prototype}, not a black-box security audit of an external deployed system.

\paragraph{Positioning.}
We explicitly avoid the term ``vulnerability disclosure'' for this work. ``Disclosure'' conventionally refers to reporting a flaw discovered in a third-party, deployed system to its maintainers. In our setting, the system is the authors' own prototype, has never been publicly released, and has no external users. The contribution is therefore better described as: \emph{(i)~we identify and characterize an authorization attack vector on the recovery sink of self-healing LLM agents (topology-mutation privilege escalation), (ii)~we instantiate and study it on a representative self-healing architecture prototype developed by the authors, and (iii)~we propose and validate a portable design pattern (\reauthgate{}) that enforces the privilege-downgrade invariant}. We frame the evidence accordingly: the validation results are measurements on the authors' own prototype under a specific configuration, not field observations from a deployed system.

\paragraph{Naming.}
The system under study is referred to throughout this paper by the placeholder \system{} (the LaTeX macro \verb|\system{}| expands to the chosen name). The name used here is a paper-level label for the prototype under study; it is not a product name and does not refer to any externally deployed system.

\subsection{Internal Remediation Timeline}

Because the attack vector is in the authors' own system, no external CVE request or third-party notification was required. The internal review and remediation timeline is:

\begin{itemize}[leftmargin=*]
    \item \textbf{2026-04-19:} V2 identified during a red-team review of the recovery pipeline. Root cause identified as the \code{parseAndValidate(validate=false)} call in \code{TopologyMutationStrategy}, which parses a diagnostic-LLM-proposed role without checking it against the privilege set authorized for the failed execution.
    \item \textbf{2026-04-26:} Defense design finalized (unified \reauthgate{} enforcing the \privinv{}). Implementation completed in a feature branch.
    \item \textbf{2026-05-10:} Defense merged to the main branch and deployed to the internal staging environment. No production incident has been attributed to V2 prior to the fix; the issue was discovered proactively, not in response to an attack.
    \item \textbf{2026-05-17:} Validation (reported in \S\ref{sec:evaluation}) started on the patched codebase; the baseline measurements were obtained by toggling the defense off via a runtime flag on the same patched code, ensuring the only difference between configurations is the defense itself.
    \item \textbf{2026-06-21:} Paper submitted for double-blind review. The patched code is the current code; no separate ``patched release'' was issued because the system has not been publicly released.
\end{itemize}

\subsection{Ethics Statement}

\paragraph{No third-party harm.} Because \system{} is the authors' own system and has not been publicly released, the experiments in this paper do not expose any third-party user to risk. No live user data was used in the validation; all attack payloads were synthetic and directed at canary tools deployed in an isolated test environment.

\paragraph{Attack payload release.} We commit to releasing the attack payloads, the validation harness, and the anonymized real-LLM response logs upon publication. We follow the principle of \emph{responsible release}: the released artifacts include only payloads that target \emph{structural} authorization properties of the recovery sink (the proposed role and its authorization set) and do not include model-specific jailbreaks that would transfer to unrelated production systems. The defense is shipped alongside the attack payloads, so that any party reproducing the attack also has the fix available.

\paragraph{LLM API usage.} The validation consumed a small number of paid API calls (approximately 90 trials across the V2 attack, benign recovery, and benign topology-mutation configurations, $N=30$ per configuration) to DeepSeek-V4-Flash (via SenseNova) and Qwen2.5-72B-Instruct (via SiliconFlow) for the \system{} testbed validation reported in \S\ref{sec:evaluation}, at a total cost of under US\$5. No model provider's terms of service were violated; the prompts used are research-purpose red-team prompts targeting our own system, not attempts to extract training data or bypass model safety training in general.

\paragraph{Communication to other framework maintainers.}
Our evidence for cross-framework impact is structural rather than exploit-based, and we calibrate our communication accordingly. The \reauthgate{} pattern applies, in principle, to any self-healing agent system that exposes a topology-mutation or role-replacement primitive on its recovery path. We frame our communication to other framework maintainers as a \emph{design recommendation} (``if your recovery path can mutate the agent topology or replace a role, consider enforcing a privilege-downgrade invariant at a centralized reauthorization gate before the side effect executes''), not as a vulnerability report on any specific third-party system. We will communicate this design recommendation, with a pointer to the pattern and reference implementation, to maintainers of widely-deployed agent frameworks at least 45 days prior to publication, and will provide the program committee with verification evidence (issue links, acknowledgment emails) upon request.

\paragraph{Broader impact.} Self-healing mechanisms are becoming a standard component of production LLM agent systems, and the recovery-path authorization boundary they introduce is currently under-recognized. By naming and characterizing this boundary, we aim to enable framework authors to design recovery pipelines with explicit privilege invariants from the outset, rather than retrofitting them after deployment. The cost of the proposed defense (sub-millisecond machinery overhead, zero false alarms) is low enough that we expect it to become a default component of future self-healing agent architectures that expose privileged topology-mutation primitives.
